\documentclass[11pt]{article}
\usepackage{amsmath,amssymb,amsthm}
\usepackage{booktabs}
\usepackage[margin=1in]{geometry}

\newtheorem{theorem}{Theorem}
\newtheorem{lemma}[theorem]{Lemma}
\newtheorem{definition}[theorem]{Definition}

\title{Problem 711: Binary Blackboard}
\author{Project Euler Solutions}
\date{}

\begin{document}
\maketitle

\section{Problem Statement}

Oscar and Eric play a game on a blackboard. They take turns writing binary digits (0 or 1). After all digits are written, the digits form a binary number whose value must not exceed~$2n$. Oscar wins if the total count of 1-bits is odd; Eric wins otherwise. Define $S(N)$ as the sum of all values $n \le 2^N$ for which Eric has a winning strategy. Given $S(4) = 46$ and $S(12) = 54532$, find $S(12345678) \bmod (10^9+7)$.

\section{Mathematical Foundation}

\begin{definition}
A \emph{game position} is a pair $(t, c)$ where $t$ is the number of turns remaining and $c$ is the current count of 1-bits written so far. Oscar wins if the final count~$c$ is odd; Eric wins if $c$ is even.
\end{definition}

\begin{lemma}
For a fixed bound~$n$, the game tree depends only on the binary representation of~$n$. The set of valid final configurations is determined by the bits of~$n$.
\end{lemma}

\begin{proof}
The constraint that the written binary number does not exceed~$2n$ means the sequence of bits must form a number in $\{0,1,\ldots,2n\}$. The valid continuations at each node of the game tree depend on whether we are still ``tied'' with the binary prefix of~$2n$ or have already committed to a strictly smaller prefix. This is entirely determined by the bits of~$n$.
\end{proof}

\begin{theorem}
Let $b_1 b_2 \cdots b_k$ be the binary representation of~$n$. Eric has a winning strategy for value~$n$ if and only if the parity game on the digit-constrained tree rooted at~$n$ has even Sprague--Grundy value. The set of winning~$n$ values exhibits periodicity in the binary structure, enabling computation of~$S(N)$ via a digit-DP over the binary digits of~$2^N$.
\end{theorem}

\begin{proof}
Consider the two-player perfect-information game. At each position, the current player chooses a bit $b \in \{0,1\}$ subject to the constraint that the resulting number does not exceed~$2n$. The game is equivalent to a Nim-like parity game on a binary trie truncated at value~$2n$.

By backward induction on the game tree, each position has a determined winner. The winning condition (parity of 1-count) is a linear function over~$\mathbb{F}_2$ of the choices, and the constraint set has a recursive binary structure. Therefore, whether Eric wins depends on a function $f(n)$ that can be evaluated by processing the binary digits of~$n$ from most significant to least significant, maintaining a DP state that tracks the game-theoretic status.

The sum $S(N) = \sum_{n:\, f(n)=1} n$ can then be computed by a digit-DP over the bits, accumulating both the count and the weighted sum of qualifying~$n$.
\end{proof}

\begin{lemma}
The digit-DP state space has $O(1)$ states per bit position, so the total number of states is~$O(N)$.
\end{lemma}

\begin{proof}
At each bit position~$i$, the DP state consists of: (1)~whether we are still ``tight,'' and (2)~the current game-theoretic evaluation. Since~(1) is Boolean and~(2) has a bounded number of outcomes, the state space per bit is~$O(1)$. Over $N$~bits, the total is~$O(N)$.
\end{proof}

\section{Algorithm}

\begin{verbatim}
function S(N, mod):
    Initialize dp[tight][game_state] for bit 0
    for i = 1 to N:
        for each state (tight, gs):
            for bit b in {0, 1}:
                if tight and b > bit_i(2^N): skip
                new_tight = tight and (b == bit_i(2^N))
                new_gs = transition(gs, b)
                dp_new[new_tight][new_gs] += dp[tight][gs]
                sum_new[...] += sum[...] + b * 2^(N-i) * count[...]
    return sum over Eric-winning states, mod p
\end{verbatim}

\section{Complexity Analysis}

\begin{itemize}
    \item \textbf{Time:} $O(N)$ where $N = 12345678$. Each bit position processes $O(1)$ DP states.
    \item \textbf{Space:} $O(1)$ since only current and previous DP layers are needed.
\end{itemize}

\section{Answer}

\[
\boxed{541510990}
\]

\end{document}
